\documentclass[12pt,a4paper,oneside]{article}
\usepackage[brazil]{babel}
\usepackage[utf8]{inputenc}
\usepackage[dvipsnames]{xcolor}
\usepackage{listings}
\usepackage{wrapfig}
\usepackage{olymp}

\setcounter{problem}{1}

\begin{document}

\begin{problem}{Contagem regressiva}{contagem}{2}

\begingroup
\setlength{\intextsep}{0pt}%
%\setlength{\columnsep}{0pt}

\begin{wrapfigure}{r}{0.4\textwidth}
    \includegraphics[width=0.4\textwidth]{figuras/artemisiicrew.png}
    \centering {\footnotesize Astronautas da missão Artemis II}
\end{wrapfigure}

Após vários adiamentos, está previsto para esta quarta-feira, 1º de abril, o lançamento do primeiro voo tripulado do Programa Artemis, da NASA.
A missão Artemis II deve partir rumo à Lua com quatro astronautas. É a primeira missão tripulada à Lua em mais de 50 anos!

Não está previsto pouso na Lua, apenas um sobrevoo ao seu redor, seguido do retorno à Terra, mas a NASA planeja transmitir a missão ao vivo e, assim, todos poderão acompanhar a contagem regressiva final do lançamento.

Faça um programa para escrever esta contagem final. Tradicionalmente, são 10 segundos de contagem final, mas seu programa deve estar preparado para ler um valor inicial qualquer e depois escrever a contagem a partir dele.

% \Notes

\InputFile

A entrada contém apenas um número natural, $N$, que indica a partida da contagem.

Restrições: $1 \leq N \leq 100$.

\endgroup

\OutputFile

Escreva uma linha na saída, com a contagem de $N$ até $0$. Os valores devem ser separados por espaço em branco, mas não deve haver espaço em branco após o $0$, apenas quebra de linha.

% \Notes

\Example

\begin{example}
\exmp{
5
}{
5 4 3 2 1 0
}%
\end{example}

\begin{example}
\exmp{
10
}{
10 9 8 7 6 5 4 3 2 1 0
}%
\end{example}

\begin{example}
\exmp{
1
}{
1 0
}%
\end{example}

%Comentários sobre o exemplo, se necessário.

\end{problem}

\end{document}